% !PRINTED 578
\weberpage{578}

%<<CONT
setzen können. Da nun in der Reihe der Zahlen $\nu^2$ die Zahl 0
und ausserdem jeder quadratische Rest zweimal vorkommt [für
$\nu^2$ und $(-\nu)^2$], so ist

\[ r^{\frac{m^2}{4}} f(r^{-2}, m) = 1 + 2A = A - B, \]

und aus der Formel (12) ergiebt sich

\[ A - B = r^{\frac{m^2}{4}} (1 - r^{-2})(1 - r^{-6}) \ldots (1 - r^{-2n + 4}). \]

Nun ist

\[ r^{\frac{m^2}{4}} = r\, r^3 r^5 \ldots r^{n - 2}, \]

und wir bekommen also schliesslich die Formel

\weq{13}{A - B = (r - r^{-1})(r^3 - r^{-3}) \ldots (r^{n - 2} - r^{-n + 2}).}

Bedeutet $\nu$ eine ungerade und $\mu$ eine gerade Zahl zwischen
0 und ${}^1/_2 n$, so besteht dies Product aus den Factoren

\[ r^{\nu} - r^{-\nu}, \quad r^{n - \mu} - r^{-n + \mu} = -(r^{\mu} - r^{-\mu}). \]

Die Anzahl der Factoren der zweiten Kategorie ist so gross,
als die Anzahl der geraden Zahlen zwischen 0 und ${}^1/_2 n$, oder
gleich der Anzahl der ganzen positiven Zahlen, die kleiner als
${}^1/_4 n$ sind. Setzen wir also, indem wir unter $l$ eine ganze Zahl
verstehen, $n = 4l + 1$ oder $4l + 3$, so ist $l$ diese Anzahl, und
wir können setzen

\[ A - B = (-1)^{l} \overset{\nu}{\Pi} (r^{\nu} - r^{-\nu}) \overset{\mu}{\Pi} (r^{\mu} - r^{-\mu}), \]

wofür auch, wenn $h$ die Reihe der Zahlen von 1 bis ${}^1/_2 (n - 1)$
durchläuft, gesetzt werden kann

\[ A - B = (-1)^{l} \overset{h}{\underset{1,\, \frac{n - 1}{2}}{\Pi}} (r^{h} - r^{-h}). \]

Setzen wir nun

\[ r = e^{\frac{2 \pi i k}{n}}, \quad r^{h} - r^{-h} = 2i \sin \frac{2 \pi k h}{n}, \]

so wird

\[ A - B = (-1)^{l} (2i)^{\frac{n - 1}{2}} \overset{h}{\Pi} \sin \frac{2 \pi h k}{n}. \]

Das Product, was hier noch vorkommt, haben wir schon
früher bestimmt [\S.\,138, (4)]:

\[ 2^{\frac{n - 1}{2}} \overset{h}{\Pi} \sin \frac{2 \pi k h}{n} = \left( \frac{k}{n} \right) \sqrt{n}, \]

worin die $\sqrt{n}$ positiv zu nehmen ist.
